
\setcounter{page}{164}
\section{Duality for a finite morphism}

Throughout this section we let \(f: X \to Y\) be a finite morphism of locally noetherian preschemes. We will define a functor \(f^b\) and a morphism of functors
\[
\operatorname{Trf}_f\colon \Rfst f^b \to \id,
\]
with the same formal properties as \(f^!\) and \(\operatorname{Trp}_f\) of \S4 above. Then we prove a duality theorem similar to the one for projective space. This duality theorem is much more elementary, but it is important to place it in the correct functorial context.

The locally noetherian hypothesis is not required merely for the definition of \(f^b\), but is essential for its functorial properties and the existence of the trace map. This suggests the definition may not be canonical in the non-noetherian setting. We also give an example showing quasi-coherence hypotheses on sheaves are necessary for duality to hold.

Let \(f: X \to Y\) be a finite morphism of locally noetherian preschemes. Let \(\overline{f}\colon (X,\sheaf{O}_X)\to (Y,f_*\sheaf{O}_X)\) be the induced morphism of ringed spaces, and let \(\operatorname{Mod}(f_*\sheaf{O}_X)\) denote the category of \(f_*\sheaf{O}_X\)-modules on \(Y\). We consider the functors
\setcounter{page}{165}
\[
\overline{f}_{*}\colon \operatorname{Mod}(X) \to \operatorname{Mod}(f_{*}\sheaf{O}_X),
\]
\[
\overline{f}^{*}\colon \operatorname{Mod}(f_{*}\sheaf{O}_X) \to \operatorname{Mod}(X).
\]

Since \(f\) is finite, \(\overline{f}\) is flat, hence \(\overline{f}^*\) is exact. The pair \((\overline{f}^*,\overline{f}_*)\) are adjoint functors [EGA O], with unit morphism
\[
\tau\colon \id \to \overline{f}_{*}\overline{f}^{*},
\]
inducing a natural isomorphism
\[
\Hom_{\sheaf{O}_X}\bigl(\overline{f}^*G, F\bigr)
\cong
\Hom_{f_*\sheaf{O}_X}\bigl(G, \overline{f}_*F\bigr)
\]
for all \(F\in\operatorname{Mod}(X)\), \(G\in\operatorname{Mod}(f_*\sheaf{O}_X)\).

\begin{definition}
Let \(f: X \to Y\) be a finite morphism of locally noetherian preschemes. Define the functor
\[
f^{b}\colon D^{+}(Y) \to D^{+}(X)
\]
by
\[
f^{b}
= \overline{f}^{*} \circ \RHom_{Y}\bigl(f_{*}\sheaf{O}_X,\,\cdot\bigr).
\]
Here \(\RHom_Y(f_*\sheaf{O}_X,\,\cdot)\) is regarded as a functor \(D^+(Y)\to D^+\bigl(\operatorname{Mod}(f_*\sheaf{O}_X)\bigr)\), and \(\overline{f}^*\) is exact so preserves derived category structure.
\end{definition}

\setcounter{page}{166}
If \(f\) has finite Tor-dimension, then \(f_*\sheaf{O}_X\) has finite Tor-dimension in \(\operatorname{Mod}(Y)\). Since \(f_*\sheaf{O}_X\) is coherent, it admits locally finite free resolutions. Consequently \(\RHom_Y(f_*\sheaf{O}_X,\,\cdot)\) has finite cohomological dimension, allowing us to extend
\[
f^{b}\colon D(Y) \to D(X)
\]
by the same formula.

\begin{proposition}
Let \(f: X \to Y\) be a finite morphism of locally noetherian preschemes (resp.\ of finite Tor-dimension). Then
\[
f^{b}\bigl(D_{\mathrm{qc}}^{+}(Y)\bigr) \subseteq D_{\mathrm{qc}}^{+}(X),\quad
f^{b}\bigl(D_{\mathrm{c}}^{+}(Y)\bigr) \subseteq D_{\mathrm{c}}^{+}(X),
\]
and in the finite Tor-dimension case
\[
f^{b}\bigl(D_{\mathrm{qc}}(Y)\bigr) \subseteq D_{\mathrm{qc}}(X),\quad
f^{b}\bigl(D_{\mathrm{c}}(Y)\bigr) \subseteq D_{\mathrm{c}}(X).
\]
\end{proposition}

\begin{proof}
Follows from [I], [II.3.2], and the fact \(\overline{f}^*\) preserves quasi-coherence and coherence of sheaves.
\end{proof}

\setcounter{page}{167}
\begin{proposition}
Let \(X\xrightarrow{f} Y\xrightarrow{g} Z\) be composable finite morphisms of locally noetherian preschemes (resp.\ finite Tor-dimension). There exists a natural functor morphism
\[
(g f)^{b} \to f^{b} g^{b}
\]
of functors \(D^+(Z)\to D^+(X)\) (resp.\ \(D(Z)\to D(X)\)). This morphism is an isomorphism on \(D_{\mathrm{qc}}^+(Z)\) (resp.\ \(D_{\mathrm{qc}}(Z)\)).
\end{proposition}

\begin{proof}
For a sheaf \(G\in\operatorname{Mod}(Z)\) we have a canonical isomorphism
\[
\overline{(gf)}^{*}\RHom_Z\bigl((gf)_*\sheaf{O}_X,G\bigr)
\cong
\overline{f}^{*}\RHom_Y\Bigl(f_*\sheaf{O}_X,\,\overline{g}^{*}\RHom_Z\bigl(g_*\sheaf{O}_Y,G\bigr)\Bigr).
\]
By [I.5.4] this induces the functor morphism \((gf)^b\to f^b g^b\).

To verify it is an isomorphism on quasi-coherent complexes, use the way-out functor lemma to reduce to injective quasi-coherent sheaves. For a ring homomorphism \(A\to B\), \(\Hom_A(B,I)\) is injective whenever \(I\) is injective over \(A\); combined with [II.7.14], [II.7.16] the isomorphism follows.
\end{proposition}

\setcounter{page}{168}
\begin{proposition}
Let \(f: X\to Y\) be finite locally noetherian, \(u: Y'\to Y\) flat locally noetherian, and form the Cartesian square
\[
\begin{CD}
X' @>v>> X \\
@VgVV @VfVV \\
Y' @>u>> Y
\end{CD}
\]
with \(X'=X\times_Y Y'\). Then there is a natural functorial isomorphism
\[
v^{*} f^{b}(G^{\cdot}) \cong g^{b} u^{*}(G^{\cdot})
\]
for all \(G^\bullet\in D^+(Y)\) (resp.\ \(G^\bullet\in D(Y)\)).
\end{proposition}

\begin{proof}
Follows from [II.5.8]. Details omitted.
\end{proposition}

\begin{corollary}
With the above hypotheses, if additionally \(u,v\) are smooth, then
\[
v^{*} f^{\dagger}(G^{\cdot}) \cong g^{b} u^{\#}(G^{\cdot}),
\]
compatible with the composition isomorphisms of Propositions 2.2 and 6.2.
\end{corollary}

\begin{proof}
Immediate from base-change compatibility of dualizing sheaves.
\end{corollary}

\setcounter{page}{169}
\begin{proposition}
Let \(f: X\to Y\) be a finite morphism of locally noetherian preschemes (resp.\ finite Tor-dimension). There exists a functorial morphism
\[
\operatorname{Trf}_{f}\colon \Rfst f^{b}(G^{\cdot}) \to G^{\cdot}
\]
for \(G^\bullet\in D_{\mathrm{qc}}^+(Y)\) (resp.\ \(G^\bullet\in D_{\mathrm{qc}}(Y)\)).
\end{proposition}

\begin{proof}
For \(G\in\operatorname{Mod}(Y)\), consider the canonical unit map
\[
\tau\colon \RHom_Y\bigl(f_*\sheaf{O}_X,G\bigr)
\to \overline{f}_{*}\overline{f}^{*}\RHom_Y\bigl(f_*\sheaf{O}_X,G\bigr).
\]
This induces a derived functor morphism
\[
\mathbf{R}\tau\colon \RHom_Y\bigl(f_*\sheaf{O}_X,G^\bullet\bigr)
\to \Rfst f^b(G^\bullet).
\]
For quasi-coherent \(G^\bullet\), \(\mathbf{R}\tau\) is an isomorphism: reduce to injective quasi-coherent sheaves via way-out lemma, using affineness of finite morphisms [EGA II.4.3].

Compose \((\mathbf{R}\tau)^{-1}\) with the evaluation-at-unit morphism
\[
\RHom_Y\bigl(f_*\sheaf{O}_X,G^\bullet\bigr) \to G^\bullet
\]
to define \(\operatorname{Trf}_f\).
\end{proposition}
